% =============================================================
% identification_template.tex
% Empirical strategy section. Lead with the identifying argument,
% then the equation, then the supporting evidence.
% =============================================================

\section{Empirical Strategy}\label{sec:strategy}

\subsection{Identification}\label{sec:identification}

Our research design exploits [SOURCE OF VARIATION: e.g.\ the staggered
rollout of program X across U.S.\ states between 2010 and 2018]. The
identifying assumption is that, absent the program, [TREATED UNITS]
would have followed the same trend in [OUTCOME] as [CONTROL UNITS].

We provide three pieces of evidence in support of this assumption.

\paragraph{Pre-trends.} We estimate the event-study specification
%
\begin{equation}\label{eq:event_study}
y_{it} = \sum_{e \neq -1} \beta_e \, \ind\{t - g_i = e\}
       + \alpha_i + \gamma_t + \varepsilon_{it},
\end{equation}
%
where $g_i$ is unit $i$'s first treatment date (with $g_i = \infty$
for never-treated units) and $\beta_e$ traces out the dynamic effect
at event time $e$. \Cref{fig:event} reports estimates from the
heterogeneity-robust Sun-Abraham (2021) implementation. Pre-period
leads are jointly indistinguishable from zero ($p = [P]$), supporting
the parallel-trends assumption.

\paragraph{Placebo.} Applying the same design to [PLACEBO OUTCOME],
which the program should not affect, yields a point estimate of
[POINT] (s.e.\ [SE]) -- a precise null (\Cref{tab:placebo}).

\paragraph{Balance.} \Cref{tab:balance} reports standardized
differences in pre-treatment characteristics across treatment cohorts.
All differences are within $\pm 0.10$, the conventional
\citet{ImbensRubin2015} threshold for ``well-balanced''.

\subsection{Estimating Equation}

Our preferred specification is
%
\begin{equation}\label{eq:main}
y_{it} = \beta \cdot D_{it} + X_{it}^{\prime} \gamma + \alpha_i + \gamma_t
       + \varepsilon_{it},
\end{equation}
%
where $D_{it}$ is an indicator for unit $i$ being treated in period $t$,
$X_{it}$ is a vector of time-varying controls, and $\alpha_i$, $\gamma_t$
are unit and time fixed effects. Standard errors are clustered at the
[CLUSTER LEVEL], the level at which treatment is assigned. With
[N CLUSTERS] clusters we also report wild cluster bootstrap p-values
(\citealp{CameronGelbachMiller2008}; \citealp{MacKinnonWebb2018}).

Because treatment timing is staggered, we estimate \cref{eq:main} using
the [Callaway-Sant'Anna (2021) / Sun-Abraham (2021) / Borusyak-Jaravel-Spiess
(2024)] heterogeneity-robust estimator and report the conventional TWFE
estimate as a benchmark in \Cref{tab:did_estimators}.

\subsection{Threats to Identification}

The remaining concern is [HONEST CONCERN: e.g.\ that policy adoption is
correlated with pre-existing trends in [CHANNEL]]. We address this in
three ways. First, [APPROACH 1]. Second, [APPROACH 2]. Third,
[APPROACH 3]. Results are reported in \Cref{sec:robustness}.
